% --- UNL TIC PIIC TEMPLATE ---
% Formato específico para Proyectos de Integración Curricular (PIIC) - UNL
\documentclass[12pt,spanish]{article}

\usepackage{graphicx}
\usepackage{geometry}
\usepackage{indentfirst} 
\usepackage{caption}
\usepackage{floatrow} 
\usepackage{booktabs}
\usepackage{float}
\usepackage{longtable}
\usepackage{threeparttable}
\usepackage{array}
\usepackage{calc}
\usepackage{amsmath, amssymb}
\usepackage{fancyhdr}
\usepackage{mathptmx} 
\usepackage[T1]{fontenc}
\usepackage{setspace} 
\usepackage{ragged2e} 
\usepackage{titlesec}
\usepackage{url}
\def\UrlBreaks{\do\/\do-\do\.\do\_\do\?\do\=\do\&\do\%}
\usepackage[hidelinks,breaklinks]{hyperref}

% --- Pandoc Compatibility Block ---
\providecommand{\tightlist}{%
  \setlength{\itemsep}{0pt}\setlength{\parskip}{0pt}}

\usepackage{xcolor}
\usepackage{fancyvrb}
\usepackage{framed}

% --- Pandoc Citeproc (CSL) compatibility ---
\newlength{\cslhangindent}
\setlength{\cslhangindent}{1.27cm} 
\newlength{\csllabelwidth}
\setlength{\csllabelwidth}{3em}
\newcommand{\citeproctext}{} 
\newcommand{\citeprocitem}[2]{\item} 
\newenvironment{CSLReferences}[2] 
 {\begin{list}{}{
  \setlength{\leftmargin}{\cslhangindent}
  \setlength{\parsep}{0pt}
  \setlength{\itemsep}{6pt}
  \ifodd #1
   \setlength{\leftmargin}{\leftmargin + \csllabelwidth}
   \setlength{\itemindent}{-\csllabelwidth}
  \fi
  }
  \renewcommand{\bibitem}[2][]{\item}
 }
 {\end{list}}
\newcommand{\CSLBlock}[1]{#1\hfill\break}
\newcommand{\CSLLeftMargin}[1]{\parbox[t]{\csllabelwidth}{#1}}
\newcommand{\CSLRightInline}[1]{\parbox[t]{\linewidth - \csllabelwidth}{#1}\break}
\newcommand{\CSLIndent}[1]{\hspace{\cslhangindent}#1}

% Solución al error 'No counter none defined' en longtable
\makeatletter
\newcounter{none}
\@ifpackageloaded{caption}{
  \DeclareCaptionType{none}
  \captionsetup[none]{name=Tabla}
}{}
\makeatother

\makeatletter
\@ifundefined{paragraph}{\newcounter{paragraph}}{}
\@ifundefined{subparagraph}{\newcounter{subparagraph}}{}
\makeatother

% --- Localización Manual ---
\AtBeginDocument{
  \renewcommand{\tablename}{Tabla}
  \renewcommand{\figurename}{Figura}
  \renewcommand{\contentsname}{Índice}
  \renewcommand{\listtablename}{Índice de Tablas}
  \renewcommand{\listfigurename}{Índice de Figuras}
  \renewcommand{\abstractname}{Resumen}
  \renewcommand{\refname}{Bibliografía} % Según lineamientos TIC
  
  \RaggedRight
  \setlength{\parindent}{1.25cm}
  \setlength{\parskip}{0pt}
}

\captionsetup{
  position=top,
  justification=RaggedRight,
  singlelinecheck=false,
  labelfont=bf,
  textfont=it,
  labelsep=newline,
  font=small,
  skip=10pt
}

% Márgenes UNL TIC
\geometry{
  a4paper,
  top=2.5cm,
  bottom=2.5cm,
  left=3.0cm,
  right=2.5cm
}

% --- Encabezados (SIN LOGOS) ---
\pagestyle{fancy}
\fancyhf{} 
\renewcommand{\headrulewidth}{0pt}
\fancyfoot[C]{\thepage} % Numeración centrada abajo según normas comunes de PIIC si no se especifica

% --- Títulos ---
\titleformat{\section}{\normalfont\normalsize\bfseries\centering}{\thesection}{1em}{}
\titleformat{\subsection}{\normalfont\normalsize\bfseries}{\thesubsection}{1em}{}
\titleformat{\subsubsection}{\normalfont\normalsize\bfseries\itshape}{\thesubsubsection}{1em}{}
\titlespacing*{\section}{0pt}{18pt}{12pt}
\titlespacing*{\subsection}{0pt}{12pt}{6pt}
\titlespacing*{\subsubsection}{0pt}{12pt}{6pt}

\makeatletter\@ifpackageloaded{caption}{}{\usepackage{caption}}\AtBeginDocument{%
\ifdefined\contentsname
  \renewcommand*\contentsname{Table of contents}
\else
  \newcommand\contentsname{Table of contents}
\fi
\ifdefined\listfigurename
  \renewcommand*\listfigurename{List of Figures}
\else
  \newcommand\listfigurename{List of Figures}
\fi
\ifdefined\listtablename
  \renewcommand*\listtablename{List of Tables}
\else
  \newcommand\listtablename{List of Tables}
\fi
\ifdefined\figurename
  \renewcommand*\figurename{Figure}
\else
  \newcommand\figurename{Figure}
\fi
\ifdefined\tablename
  \renewcommand*\tablename{Table}
\else
  \newcommand\tablename{Table}
\fi
}\@ifpackageloaded{float}{}{\usepackage{float}}
\floatstyle{ruled}
\@ifundefined{c@chapter}{\newfloat{codelisting}{h}{lop}}{\newfloat{codelisting}{h}{lop}[chapter]}
\floatname{codelisting}{Listing}\newcommand*\listoflistings{\listof{codelisting}{List of Listings}}\makeatother\makeatletter\makeatother\makeatletter\@ifpackageloaded{caption}{}{\usepackage{caption}}
\@ifpackageloaded{subcaption}{}{\usepackage{subcaption}}\makeatother

\begin{document}

% --- Portada PIIC ---
\begin{titlepage}
    \begin{center}
        \setstretch{1.0}
        {\bfseries \Large UNIVERSIDAD NACIONAL DE LOJA}\\[0.5cm]
        {\bfseries \large FACULTAD JURÍDICA, SOCIAL Y ADMINISTRATIVA}\\[0.3cm]
        {\bfseries \large CARRERA DE ECONOMÍA}\\[2.5cm]
        
        {\bfseries \large PROYECTO DE INVESTIGACIÓN DE INTEGRACIÓN CURRICULAR}\\[1.5cm]
        
        \setstretch{1.5}
        {\bfseries \Large Informe Metodológico: Desigualdad y
Concentración del Ingreso en Pichincha (2025)}\\[2.5cm]
    \end{center}
    
    \setstretch{1.2}
    \vfill
    \noindent {\bfseries Autor:} Erick Fabricio Condoy Seraquive \\
    \noindent {\bfseries Tutor:} Econ. Jose Rafael Alvarado Lopez \\
    
    \vspace{2cm}
    
    \begin{center}
        {\bfseries LOJA -- ECUADOR}\\[0.2cm]
        {\bfseries 2026}
    \end{center}
\end{titlepage}

\newpage
\setstretch{1.5}
\pagenumbering{arabic}
\setcounter{page}{1}

\section{Apartado 1: Medición de la Desigualdad mediante el Coeficiente
de
Gini}\label{apartado-1-mediciuxf3n-de-la-desigualdad-mediante-el-coeficiente-de-gini}

La medición de la desigualdad económica es un componente crítico para
comprender los determinantes de la doble informalidad en el Ecuador.
Para la provincia de \textbf{Pichincha}, utilizando la base anualizada y
acumulada de la \textbf{ENEMDU 2025}, estimamos la concentración del
ingreso per cápita del hogar.

\subsection{Metodología y Formulación
Matemática}\label{metodologuxeda-y-formulaciuxf3n-matemuxe1tica}

El Coeficiente de Gini se calcula a partir de los microdatos expandidos
utilizando la formulación ponderada de Brown. Sea
\(Y = (y_1, y_2, \dots, y_n)\) el vector de ingresos ordenado de forma
ascendente (\(y_1 \le y_2 \le \dots \le y_n\)) y
\(W = (w_1, w_2, \dots, w_n)\) el vector de factores de expansión
asociados (ponderaciones poblacionales).

Definimos las proporciones acumuladas de población (\(p_i\)) e ingresos
(\(q_i\)) para cada observación \(i\) como:

\[p_i = \frac{\sum_{j=1}^{i} w_j}{\sum_{j=1}^{N} w_j}\]

\[q_i = \frac{\sum_{j=1}^{i} w_j y_j}{\sum_{j=1}^{N} w_j y_j}\]

El Coeficiente de Gini ponderado está dado por la relación:

\[\text{Gini} = 1 - \sum_{i=1}^{N} f_i \left( q_i + q_{i-1} \right)\]

Donde \(f_i = p_i - p_{i-1} = \frac{w_i}{\sum_{j=1}^{N} w_j}\)
representa la fracción de población que aporta la observación \(i\), con
\(q_0 = 0\).

\subsection{Tratamiento y Sanitización de
Microdatos}\label{tratamiento-y-sanitizaciuxf3n-de-microdatos}

Para garantizar estimaciones insesgadas y consistentes con los
estándares del Instituto Nacional de Estadística y Censos (INEC),
aplicamos el siguiente protocolo en la base de datos:

\begin{enumerate}
\def\labelenumi{\arabic{enumi}.}
\tightlist
\item
  \textbf{Filtro Geográfico:} Se seleccionó a la población residente en
  la provincia de Pichincha filtrando por la variable de localización
  \texttt{ciudad} (código de provincia \texttt{17}).
\item
  \textbf{Tratamiento de Valores Nulos:} Se eliminaron de forma estricta
  las observaciones con valores perdidos (\texttt{missing}) en la
  variable de ingreso per cápita del hogar (\texttt{ingpc}).
\item
  \textbf{Depuración de Ingresos Negativos:} Se removieron los valores
  menores a cero (\texttt{ingpc\ \textless{}\ 0}) para evitar
  distorsiones matemáticas en los denominadores acumulativos de la
  curva.
\item
  \textbf{Representatividad Muestral:} La muestra final sanitizada
  consta de \textbf{46,728 registros válidos}, con una población
  expandida total de \textbf{3'375,369.56 habitantes}.
\end{enumerate}

\subsection{Resultados del Coeficiente de
Gini}\label{resultados-del-coeficiente-de-gini}

La estimación arrojó un \textbf{Coeficiente de Gini de 0.4532}. Este
resultado refleja una moderada concentración del ingreso en la provincia
de Pichincha y sirve como línea base para los modelos econométricos
multivariantes.

\begin{center}\rule{0.5\linewidth}{0.5pt}\end{center}

\section{Apartado 2: Construcción y Visualización de la Curva de
Lorenz}\label{apartado-2-construcciuxf3n-y-visualizaciuxf3n-de-la-curva-de-lorenz}

La Curva de Lorenz es la representación gráfica de la función de
distribución acumulada de ingresos. Permite contrastar visualmente la
distribución observada frente a una distribución hipotética de igualdad
perfecta.

\subsection{Puntos de la Curva de
Lorenz}\label{puntos-de-la-curva-de-lorenz}

Los puntos acumulados \((p_i, q_i)\) representan la proporción de
ingresos acumulados obtenidos por el percentil de población
correspondiente. La diagonal de \(45^\circ\) que conecta los puntos
\((0,0)\) y \((1,1)\) representa la \textbf{Línea de Igualdad Perfecta}
(donde el \(x\%\) de la población percibe exactamente el \(x\%\) del
ingreso total).

\subsection{Consistencia y Validación
Interplataforma}\label{consistencia-y-validaciuxf3n-interplataforma}

Para asegurar la reproducibilidad y robustez de los cálculos en el
pipeline institucional, se generaron las visualizaciones de la curva de
Lorenz de manera independiente en \textbf{Python} (a través de
\texttt{matplotlib} y cargando datos en formato \texttt{Parquet}) y en
\textbf{Stata} (a través de \texttt{twoway} y cargando datos en formato
\texttt{DTA}). Ambos análisis producen curvas idénticas con un
coeficiente de Gini coincidente a 4 decimales (\textbf{0.4532}).

\begin{figure}

\begin{minipage}[t]{0.50\linewidth}

\begin{figure}[H]

\centering{

\includegraphics[width=0.9\linewidth,height=\textheight,keepaspectratio]{assets/curva_lorenz_2025.png}

}

\caption{\label{fig-python}Curva de Lorenz generada en Python}

\end{figure}%

\end{minipage}%
%
\begin{minipage}[t]{0.50\linewidth}

\begin{figure}[H]

\centering{

\includegraphics[width=0.9\linewidth,height=\textheight,keepaspectratio]{assets/curva_lorenz_2025_stata.png}

}

\caption{\label{fig-stata}Curva de Lorenz generada en Stata}

\end{figure}%

\end{minipage}%

\end{figure}%

Ambas gráficas confirman una distribución coherente donde el área de
concentración (el espacio entre la diagonal de igualdad perfecta y la
curva observada) es idéntica en ambas plataformas analíticas.

\end{document}
